Benchmark contamination is widely assumed to inflate language model scores, yet its actual effect has never been measured across model families. We perform controlled supervised fine-tuning (SFT) injection at five dosage levels into six models (7B to 14.7B) from five open-source families and discover a spectrum of dose-response patterns, not uniform inflation. Two patterns replicate across independent families: Immune (Phi-4, LLaMA-3.1) and Benefit (statistically significant for Gemma-2, bootstrap 95\% CI [$+$1.49, $+$3.14]\,pp; exploratory for OLMo-2 and Falcon3). Collapse and V-recovery are each observed in one family (Qwen2.5) across three seeds but remain preliminary pending independent replication. Effects are benchmark-conditional: MMLU-targeted contamination improves GSM8K in four of five non-excluded families, even those harmed on MMLU. Difference-in-differences analysis detects no item-specific memorization; a volume-matched control attributes 45\% of Benefit to data quantity. ContamCal, our diagnostic framework combining multi-signal exposure mapping with contamination-augmented Item Response Theory (IRT), enables within-family domain-level contamination diagnosis with credible intervals. On Qwen2.5, Social Sciences contamination effects are negative while STEM effects are positive, a pattern invisible to aggregate scoring. Out-of-sample validation on LLaMA-3.1-8B confirms that model family characteristics beyond baseline ability modulate vulnerability. Model-specific dose-response profiling is necessary for reliable benchmark interpretation\footnote{Code and data are available at \url{https://github.com/anonymous-nlp-2026-2/contamcal-irt-calibration}.}.
